\section{Proof}
\label{appendix-proof}
\begin{lemma}
  \label{appendix-transition-safety-lemma}
  Given a system $\mathfrak{S} = (\mathcal{X}, \mathcal{X}_0, f)$, a trajectory $\tau=\langle x_0, x_1, \dots \rangle \in \mathfrak{T}(\mathfrak{S})$, a labelling function $\mathfrak{L}:\mathcal{X} \rightarrow \Sigma$ and an NBA $\mathcal{A}=(\Sigma, Q, q_0, \delta, Acc)$, if a transition safety certificate $B_{k}(x, q)$ with respect to $q_k \in Acc^{start}$ can be found, then it can be guaranteed that all runs of all \emph{words} in $\mathfrak{L}(\mathfrak{S})$ cannot visit the state $q_k$.
\end{lemma}

\begin{proof}
  By Definition~\ref{transition-safety-certificate-definition}, we have $\mathcal{Y}_u = \{(x, q_k) \mid x \in \mathcal{X}\}$ as the unsafe set in the product system $\mathfrak{S} \times \mathcal{A}$. We verify that $B_k$ satisfies the barrier certificate conditions (\ref{bsinn0}, \ref{bsun0}, \ref{bstnn0}) with respect to $\mathcal{Y}_u$:
  \begin{itemize}
    \item Condition (\ref{btsc1}) states $B_k(x_0, q_0) \geq 0$ for all $(x_0, q_0) \in \mathcal{Y}_0$, which corresponds to condition (\ref{bsinn0}).
    \item Condition (\ref{btsc2}) states $B_k(x, q_k) < 0$ for all $(x, q_k) \in \mathcal{Y}_u$, which corresponds to condition (\ref{bsun0}).
    \item Condition (\ref{btsc3}) states $B_k(x, q) \geq 0 \Rightarrow B_k(x', q') \geq 0$ for all $(x', q') \in F(x, q)$, which corresponds to condition (\ref{bstnn0}).
  \end{itemize}

  By Theorem~(\ref{corollary_tbsc}), any trajectory of the product system $\mathfrak{S} \times \mathcal{A}$ cannot visit $\mathcal{Y}_u$.

  Now we establish the connection to automaton runs. For any word $w = \mathfrak{L}(\tau) \in \mathfrak{L}(\mathfrak{S})$ where $\tau = \langle x_0, x_1, \dots \rangle$ is a trajectory of $\mathfrak{S}$, the corresponding run $r$ of $w$ on $\mathcal{A}$ is a sequence $r = \langle (r_0, w_0, r_1), (r_1, w_1, r_2), \ldots \rangle$ where $r_0 = q_0$. The sequence of product states $\langle (x_0, r_0), (x_1, r_1), \ldots \rangle$ forms a trajectory in $\mathfrak{S} \times \mathcal{A}$. Since this trajectory cannot visit $\mathcal{Y}_u = \{(x, q_k) \mid x \in \mathcal{X}\}$, we have $r_i \neq q_k$ for all $i$. Therefore, runs of all words in $\mathfrak{L}(\mathfrak{S})$ cannot visit state $q_k$.
\end{proof}

\begin{lemma}
  \label{appendix-transition-persistence-lemma}
  Given a system $\mathfrak{S}=(\mathcal{X}, \mathcal{X}_0, f)$, an NBA $\mathcal{A} = (\Sigma, Q, q_0, \delta, Acc)$, and $(k, l) \in Acc^{pair}$, if a transition persistence certificate $B_{k, l}$ can be found, then it can be guaranteed that all runs of the words in $\mathfrak{L}(\mathfrak{S})$ make accepting transitions from $q_k$ to $q_l$ happen finitely often.
\end{lemma}

\begin{proof}
We define $\mathcal{X}_{k, l}=\{x \mid (q_k, \mathfrak{L}(x), q_l) \in Acc^{k,l}\}$ as the set of states where accepting transitions from $q_k$ to $q_l$ can occur.

Consider any trajectory $\tau = \langle x_0, x_1, \dots \rangle \in \mathfrak{T}(\mathfrak{S})$ and its corresponding word $w = \mathfrak{L}(\tau)$. Let $r$ be a run of $w$ on $\mathcal{A}$ starting from $q_0$. An accepting transition from $q_k$ to $q_l$ can occur at step $i$ only if $(q_k, \mathfrak{L}(x_i), q_l) \in Acc^{k,l}$, which is equivalent to $x_i \in \mathcal{X}_{k,l}$. Therefore, the number of accepting transitions from $q_k$ to $q_l$ is bounded by the number of times the trajectory visits $\mathcal{X}_{k,l}$.

By Definition~\ref{transition-persistence-certificate-definition}, we have:
\begin{itemize}
  \item $B_{k,l}(x_0) \geq 0$ by condition (\ref{btpc1}).
  \item $B_{k,l}(x_i) \geq 0$ for all $i$ by condition (\ref{btpc2}), since $B_{k,l}$ remains non-negative along the trajectory.
  \item $B_{k,l}(x_i) \geq B_{k,l}(x_{i+1}) + I_{k,l}(x_i)\epsilon$ by condition (\ref{btpc3}), where $I_{k,l}(x_i) = 1$ if $x_i \in \mathcal{X}_{k,l}$ and $0$ otherwise.
\end{itemize}

Suppose the trajectory visits $\mathcal{X}_{k,l}$ at indices $i_1, i_2, \dots, i_m$ where $m$ is the number of times accepting transitions occur. At each visit, $B_{k,l}$ decreases by at least $\epsilon$:
\begin{align}
B_{k,l}(x_0) \geq B_{k,l}(x_{i_1}) + \epsilon \geq B_{k,l}(x_{i_2}) + 2\epsilon \geq \cdots \geq B_{k,l}(x_{i_m}) + m\epsilon.
\end{align}

Since $B_{k,l}(x_{i_m}) \geq 0$ by condition (\ref{btpc2}), we have $B_{k,l}(x_0) \geq m\epsilon$, which gives $m \leq \frac{B_{k,l}(x_0)}{\epsilon}$. Therefore, the number of accepting transitions from $q_k$ to $q_l$ is bounded, ensuring they occur only finitely often.
\end{proof}


\section{Experiment}
\label{appendix-experiment}
\subsection{Template for Two-Dimensional Kuramoto Oscillator System}
The transition safety certificate template for the two-dimensional Kuramoto Oscillator system is given by:
\begin{align*}
B_0(x_1, x_2, q) = &\, c_0 + c_1 \cdot x_1 + c_2 \cdot x_2^2 + c_3 \cdot x_1 x_2 + c_4 \cdot x_1^3 \\
                   &+ c_5 \cdot q + c_6 \cdot q^2 + c_7 \cdot q x_1^2 + c_8 \cdot x_2 q
\end{align*}

\subsection{Template for Two-Room Temperature System}
For the \textit{transition safety certificate}:
\begin{align*}
B_1(x_1, x_2, q) = &\, c_0 + c_1 \cdot I_{X_0}(x_1, x_2) + c_2 \cdot I_{X_{VF}}(x_1, x_2) \\
                   &+ c_3 \cdot x_1 I_{X_0}(x_1, x_2) + c_4 \cdot x_2 I_{X_0}(x_1, x_2) \\
                   &+ c_5 \cdot x_1 I_{X_{VF}}(x_1, x_2) + c_6 \cdot x_2 I_{X_{VF}}(x_1, x_2) \\
                   &+ c_7 \cdot \max(x_1, x_2) + c_8 \cdot x_1^2 + c_9 \cdot x_2^2 + c_{10} \cdot q + c_{11} \cdot q^2.
\end{align*}
The template for the \textit{transition persistence certificate} is as follows:
\begin{align*}
B_{1,1}(x_1, x_2) = & \, c_0 + c_1 \cdot I_{X_0}(x_1, x_2) + c_2 \cdot I_{X_{VF}}(x_1, x_2) \\
                    & + c_3 \cdot x_1 I_{X_0}(x_1, x_2)) + c_4 \cdot x_2 I_{X_0}(x_1, x_2) \\
                   	& + c_5 \cdot x_1 I_{X_{VF}}(x_1, x_2) + c_6 \cdot x_2 I_{X_{VF}}(x_1, x_2) \\ 
					& + c_7 \cdot \max(x_1, x_2) + c_8 \cdot x_1^2 + c_9 \cdot x_2^2.
\end{align*}

\subsection{Neural Network Certificate Synthesis Results}
\label{app:neural_cert_results}

This section provides the neural network certificates synthesized in the three case studies.

	\subsubsection{Kuramoto Oscillator: One-Dimensional System}
	\label{app:ko1d_neural}
	\textbf{Valid transition safety certificate network parameters:}
	\begin{itemize}
	    \item \textbf{Network architecture:} 2 -> 10 -> 1 (ReLU activation network)
	    \item \textbf{Layer 1 weights} $W^{(1)} \in \mathbb{R}^{10 \times 2}$:
	    \[
	    W^{(1)} = \begin{bmatrix}
	     0.36732605 & -0.29466873 \\
	     0.44511154 & -0.10065039 \\
	    -0.2463612  &  0.13090049 \\
	     0.5710599  &  0.18457386 \\
	    -0.16401595 &  0.800196   \\
	     0.16680917 & -0.54196787 \\
	     0.03139566 & -0.581802   \\
	    -0.2696909  & -0.6400563  \\
	     0.6605623  &  0.6877757  \\
	    -0.12304431 &  0.4023705
	    \end{bmatrix}
	    \]
	    
	    \item \textbf{Layer 1 bias} $b^{(1)} \in \mathbb{R}^{10}$:
	    \[
	    b^{(1)} = \begin{bmatrix}
	     0.3557 & -0.1530 & -0.2855 & -0.3364 &  0.1495 \\
	     0.1568 &  0.2077 &  0.4496 & -0.1997 & -0.8233
	    \end{bmatrix}^\top
	    \]
	    
	    \item \textbf{Layer 2 weights} $W^{(2)} \in \mathbb{R}^{1 \times 10}$:
	    \[
	    W^{(2)} = \begin{bmatrix}
	     0.12380187 &  0.04241462 & -0.10408211 & -0.1692231  &  0.32869837 \\
	    -0.3943727  & -0.09415235 & -0.1155376  &  0.12839615 & -0.03449793
	    \end{bmatrix}
	    \]
	    
	    \item \textbf{Layer 2 bias} $b^{(2)} \in \mathbb{R}$:
	    \[
	    b^{(2)} = -0.12770885
	    \]
	    
	    \item \textbf{Analytical expression:}
	    \[
	    B(x, q) = W^{(2)} \cdot \text{ReLU}(W^{(1)} \cdot [x, q]^\top + b^{(1)}) + b^{(2)}
	    \]
	\end{itemize}

	\subsubsection{Kuramoto Oscillator: Two-Dimensional System}
	\label{app:ko2d_neural}
	
	\textbf{Valid transition safety certificate network parameters:}
	
	\begin{itemize}
	    \item \textbf{Network architecture:} 3 -> 15 -> 1 (ReLU activation network)
	    \item \textbf{Layer 1 weights} $W^{(1)} \in \mathbb{R}^{15 \times 3}$:
	    {\footnotesize
	    \[
	    W^{(1)} = \begin{bmatrix}
	     0.16947323 & -0.21574873 &  0.13164099 \\
	    -0.08310825 & -0.01274627 & -0.04113036 \\
	    -0.06397208 & -0.54293144 & -0.02944631 \\
	    -0.31079152 &  0.06346325 & -0.08720846 \\
	    -0.08576902 &  0.09323277 &  0.05105818 \\
	     0.3797034  & -0.6118731  &  0.02912938 \\
	    -0.12490943 & -0.05650124 & -0.19199097 \\
	    -1.1089448  &  0.63800216 & -0.06427222 \\
	    -0.17578183 &  0.6446198  & -0.03190435 \\
	     0.14327501 &  0.08219448 & -0.22752807 \\
	    -0.08726892 &  0.63300693 &  0.099369   \\
	    -0.01736208 &  0.02708694 &  0.01039273 \\
	     0.77797484 & -0.3705437  & -0.0628436  \\
	     0.49392858 &  0.17566358 &  0.6906795  \\
	    -0.65604967 & -0.7597173  &  0.65568864
	    \end{bmatrix}
	    \]}
	    
	    \item \textbf{Layer 1 bias} $b^{(1)} \in \mathbb{R}^{15}$:
	    {\scriptsize
	    \[
	    b^{(1)} = \begin{bmatrix}
	    -0.3938 & -0.0656 & -0.0435 & -0.3020 & -0.1906 & -0.0847 & -0.0903 \\
	    -0.3402 & -1.0316 & -0.5757 & -1.0281 & -0.0510 & -0.7655 & -0.6736 \\
	     0.2616
	    \end{bmatrix}^\top
	    \]}
	    
	    \item \textbf{Layer 2 weights} $W^{(2)} \in \mathbb{R}^{1 \times 15}$:
	    {\footnotesize
	    \[
	    W^{(2)} = \begin{bmatrix}
	     0.00767429 & -0.01955141 & -0.21718296 &  0.17923115 &  0.00531028 \\
	    -0.35746995 &  0.00876291 & -1.0855321  & -1.4935373  & -0.13878155 \\
	    -0.44775316 & -0.00077066 & -0.72719783 &  0.2176193  &  0.21799223
	    \end{bmatrix}
	    \]}
	    
	    \item \textbf{Layer 2 bias} $b^{(2)} \in \mathbb{R}$:
	    \[
	    b^{(2)} = -0.14020069
	    \]
	    
	    \item \textbf{Analytical expression:}
	    \[
	    B(x_1, x_2, q) = W^{(2)} \cdot \text{ReLU}(W^{(1)} \cdot [x_1, x_2, q]^\top + b^{(1)}) + b^{(2)}
	    \]
	\end{itemize}

\subsubsection{Two-Room Temperature System}
\label{app:two_room_neural}
\textbf{Valid transition persistence certificate network parameters:}

\begin{itemize}
    \item \textbf{Network architecture:} 2 -> 3 -> 1
    \item \textbf{Decrease constant $\epsilon$:} 0.01 (fixed)
    \item \textbf{Layer 1 weights} $W^{(1)} \in \mathbb{R}^{3 \times 2}$:
    \[
    W^{(1)} = \begin{bmatrix}
     0.09350396 & -0.06689734 \\
     0.3247081  & -0.3753044  \\
     0.8092528  & -0.7572314
    \end{bmatrix}
    \]
    
    \item \textbf{Layer 1 bias} $b^{(1)} \in \mathbb{R}^{3}$:
    \[
    b^{(1)} = \begin{bmatrix}
    -0.76164836 & 1.549208 & -1.5467515
    \end{bmatrix}^\top
    \]
\end{itemize}

\subsection{Comparative Method Setting}
\subsubsection{Closure Certificates Method}
The synthesis process uses Z3 to find candidate certificates and dReal to find counterexamples, with a dReal precision of 0.0001.

\textbf{One-Dimensional Kuramoto Oscillator System:}
The NBA is defined as $(\Sigma, Q, q_0, \delta, \text{Acc})$ to represent $F a$.
\begin{align*}
    \text{Acc} &= \{0\} \\
    Q &= \{1, 0\} \\
    q_0 &= 1 \\
    \Sigma &= \{\emptyset, \{a\}\} \\
    \delta &= \{(1, \emptyset, 1), (1, \{a\}, 0), (0, \emptyset, 0), (0, \{a\}, 0)\} \\
	a &\text{ is defined as } x \in X_{\text{u}}
\end{align*}
The certificate templates are:
\begin{align*}
    T_{0,0} &= c_0 + c_1 x + c_2 y \\
    T_{0,1} &= c_3 + c_4 x + c_5 y \\
    T_{1,0} &= c_6 + c_7 x + c_8 y \\
    T_{1,1} &= c_9 + c_{10} x + c_{11} y
\end{align*}
The final synthesized coefficients are shown in Table~\ref{tab:coeff_1d}.

\begin{table}[h]
\centering
\caption{Coefficients for the synthesized closure certificates (One-Dimensional System)}
\label{tab:coeff_1d}
\begin{tabular}{lccc}
\hline
\textbf{Template} & $c_0$ & $c_1$ & $c_2$ \\
\hline
$T_{0,0}(x, y)$ & 0.000000 & 0.000000 & 0.000000 \\
$T_{0,1}(x, y)$ & -0.500000 & 0.000000 & 0.000000 \\
$T_{1,0}(x, y)$ & -4.310550 & 1.764117 & 0.000000 \\
$T_{1,1}(x, y)$ & 0.000000 & 0.000000 & 0.000000 \\
\hline
\end{tabular}
\end{table}

\textbf{Two-Dimensional Kuramoto Oscillator System:}
The certificate template is defined for $i, j \in \{0, 1\}$ as:
\begin{align*}
T_{i,j}((x_1, x_2), (y_1, y_2)) = &\, c_{0, i, j} + c_{1, i, j} y_1 I_{X_0}(x_1, x_2) + c_{2, i, j} y_2 I_{X_0}(x_1, x_2) \\
                                  &+ c_{3, i, j} y_1 I_{X_u}(x_1, x_2) + c_{4, i, j} y_2 I_{X_u}(x_1, x_2) \\
                                  &+ c_{5,i,j} y_1 + c_{6,i,j} y_2
\end{align*}

\textbf{Two-Room Temperature System:}
The NBA for $a_0 \land F a_1$ is $(\Sigma, Q, q_0, \delta, \text{Acc})$.

\begin{align*}
    Q &= \{0,1,2,3\} \\
    \Sigma &= \{\emptyset, \{a_0\}, \{a_1\}, \{a_0, a_1\}\} \\
    q_0 &= 0 \\
    \text{Acc} &= \{2\} \\
    \delta &= \{(q_0, \{a_0\}, q_1), (q_0, \{a_1\}, q_1), (q_0, \{a_0, a_1\}, q_2), (q_0, \emptyset, q_3), \\
    &\quad\;(q_0, \{a_1\}, q_3), (q_1, \emptyset, q_1), (q_1, \{a_0\}, q_1), (q_1, \{a_1\}, q_2), \\
    &\quad\;(q_1, \{a_0, a_1\}, q_2), (q_2, \emptyset, q_2), (q_2, \{a_0\}, q_2), (q_2, \{a_1\}, q_2), \\
    &\quad\;(q_2, \{a_0, a_1\}, q_2), (q_3, \emptyset, q_3), (q_3,\{a_0\},q_3), (q_3, \{a_1\}, q_3), \\
    &\quad\;(q_3, \{a_0, a_1\}, q_3)\} \\
	a_0 &\text{ is defined as } x \in X_{0} \\
	a_1 &\text{ is defined as } x \in X_{VF} \\
\end{align*}
The certificate template for this system is defined for $i, j \in \{0, 1, 2, 3\}$ as:
\begin{align*}
T_{i,j}((x_1, x_2), (y_1, y_2)) = &\, c_{0, i, j} + c_{1, i, j}x_1 + c_{2, i, j}x_2 + c_{3, i, j}y_1 + c_{4, i, j}y_2 \\
                                  &+ c_{5, i, j}\max(x_1, x_2) + c_{6, i, j}\max(y_1, y_2) \\
                                  &+ c_{7, i, j}x_1^2 + c_{8, i, j}x_2^2 + c_{9, i, j}y_1^2 + c_{10, i, j}y_2^2
\end{align*}

\subsubsection{Neural Closure Certificates Method}
The configuration for the neural closure certificates is summarized in Table~\ref{tab:ncc_config}. A sampling density $\xi$ of 0.01 and a safety margin $\eta$ of 0.01 are used for all systems. The Lipschitz constant is estimated using the spectral norm-based algorithm.

\begin{table}[h]
\centering
\caption{Neural network architecture and parameters for Neural Closure Certificates}
\label{tab:ncc_config}
\begin{tabular}{lccc}
\hline
\textbf{System} & \textbf{Network Architecture} & \textbf{Activation} & \textbf{Lipschitz Estimation} \\
\hline
1D Kuramoto Oscillator & (4-10-1) & ReLU & Spectral Norm \\
2D Kuramoto Oscillator & (6-10-1) & ReLU & Spectral Norm \\
Two-Room Temperature & (6-10-1) & ReLU & Spectral Norm \\
\hline
\end{tabular}
\end{table}

